
% -----------------------------------------------------------
\section{Constraints Without Reasons}

Institutional drift is a familiar observation. Organizations formed around
a clear purpose accumulate procedures, departments, metrics, and requirements
until the original purpose becomes difficult to locate inside the structure
that was built to serve it. Universities formed to pursue knowledge optimize
publication counts. Regulatory agencies formed to protect the public expand
procedural compliance. Professional bodies formed to maintain standards develop
elaborate certification requirements disconnected from the competences those
requirements were meant to ensure. Companies formed to solve a concrete
problem maintain internal reporting structures long after the problem has
changed or disappeared.

The standard explanations for this pattern appeal to agents. Capture: those
who administer rules develop interests in perpetuating them. Bureaucracy:
formal structures generate their own momentum independent of the goals they
ostensibly serve. Bad incentives: the metrics available for measuring
institutional performance diverge from the objectives the institution was
formed to pursue. Agency problems: the people who run institutions are not
identical with the people the institutions were formed to benefit, and their
interests differ. These explanations are not wrong. Capture happens. Bureaucracy
is real. Incentive misalignment is common.

But these explanations are incomplete in a way that matters. They treat
institutional drift as contingent --- as something that happens when agents
behave badly or when incentive structures are poorly designed. They imply
that better agents, better incentives, or better monitoring would prevent
the failure. What they do not explain is why institutional drift is so
universal, so persistent, and so difficult to reverse even in institutions
that begin with good agents, good incentives, and serious monitoring. They
do not explain why the failure so often occurs invisibly, without any point
at which someone decided to abandon the original objective.

The structural explanation begins differently. Constraints are durable.
Reasons are fragile. A rule can be copied into a new document, transmitted
to a new generation of administrators, and enforced by people who have never
encountered the situation the rule was designed to address. What cannot be
copied by transcription alone is the argument that made the rule intelligible
--- the understanding of why this constraint, in this domain, was necessary.
That understanding requires not only the words of the original justification
but the background knowledge, practical experience, and interpretive capacity
required to make those words meaningful.

Institutional drift, on this account, is not primarily a failure of agents.
It is a failure of transmission. What decays is not the constraint but its
reason. And once the reason is gone, the constraint is available for a new
purpose: not the purpose it was designed to serve, but whatever purpose the
current administrators find locally convenient to attribute to it.

The asymmetry in transmission cost explains why this failure is structural
rather than contingent. Copying a constraint costs approximately nothing.
Every generation can inherit the rule by transcription. Transmitting the
reason for a constraint costs far more: it requires extended engagement, a
recipient capable of understanding the original problem, and a context rich
enough to make the reasoning legible. Each generation must re-earn the reason
through something closer to reconstruction than inheritance. Constraint
transmission is mechanical. Reason transmission is social. Under any pressure
that reduces the time and attention available for the social work of
transmission, reasons decay while constraints persist. The institution retains
its shape while losing its orientation.

% -----------------------------------------------------------
\section{Search Velocity and Contextual Depth}

The previous essay in this pair established that search and orientation scale
at different rates. When the cost of generating candidates falls, the set of
reachable states expands, but the capacity to evaluate trajectories does not
expand proportionally. The bottleneck shifts from production to navigation.

The same asymmetry appears at institutional scale, but its form is different.
The relevant scarcity is not between generating options and choosing among them.
It is between the accumulation of institutional decisions and the preservation
of the context that made those decisions intelligible. Institutions generate
constraints continuously --- through policies, precedents, procedures, commitments,
and agreements. What they accumulate more slowly, and lose more easily, is
the body of reasons that made each constraint appropriate in its original context.

It is useful to characterize different institutional forms along a single
axis: the ratio of search velocity to contextual depth. Search velocity
measures how rapidly an institution can produce and implement new actions,
decisions, and policies. Contextual depth measures how much accumulated
reasoning must be consulted before an action is considered legitimate.

\begin{principle}[Contextual Depth]
Institutional responsiveness increases as required contextual depth decreases,
but orientation becomes increasingly vulnerable to projection collapse.
Institutions differ not only in efficiency but in how much of their accumulated
orientation must be engaged before legitimate action can be taken.
\end{principle}

High search velocity with low contextual depth enables rapid adaptation but
provides little protection against projection collapse. Each new decision is
made in a thin context. Early constraints are quickly overwritten rather than
built upon. The institution remains agile but its trajectory is easily deflected
by local conditions. Low search velocity with high contextual depth provides
strong protection against projection collapse but may prevent the institution
from responding to genuine changes in its situation. Each new decision must
negotiate a dense accumulated context that may itself have drifted.

This is a pressure, not a law. Institutions can resist it. But resisting it
requires deliberate investment in exactly the mechanisms that search-velocity
pressure tends to erode: the maintenance of interpretive communities who can
engage the accumulated context, the preservation of reasons alongside rules,
and the cultivation of trajectory assessment as a distinct institutional practice.
Most institutions do not make this investment systematically, because the
return is invisible until the capacity is lost.

% -----------------------------------------------------------
\section{The Reason Retention Principle}

The structural analysis suggests a precise criterion for institutional
orientation.

\begin{principle}[Reason Retention]
An institution remains oriented only to the extent that it can reconstruct
the arguments that originally justified its constraints.
\end{principle}

The word \emph{reconstruct} is doing essential work here and must be
distinguished from two weaker capacities that are often confused with it.

The first weaker capacity is \emph{archival retention}: the preservation of
documents, records, and stated rules. A sufficiently large archive contains
the original justifications in some form. But archival retention does not
guarantee that anyone can read those justifications in context, understand
the problem they were responding to, or recognize their relevance to current
decisions. Archives preserve words. Reconstruction requires understanding.

The second weaker capacity is \emph{recitation}: the ability to reproduce
the stated rationale for a rule on demand. Members of an institution may
be able to state why a procedure exists --- "this requirement ensures quality
control" or "this review process protects against conflict of interest" ---
without having the interpretive capacity to apply that rationale to novel
situations, to recognize when the rationale no longer applies, or to
distinguish cases where the original constraint is genuinely relevant from
cases where it has become ritual. Recitation is the form without the capacity.

Reconstruction in the full sense requires something stronger: the ability
to regenerate the original argument from first principles, to identify the
problem the constraint was designed to address, to understand why other
approaches were considered inadequate, and to apply that understanding to
situations the original designers did not anticipate. An institution that
can reconstruct in this sense can determine when a constraint should be
preserved, when it should be revised, and when it has outlived its original
justification. An institution that can only recite cannot make these
determinations reliably.

\begin{observation}
An institution may lose reconstructive capacity while retaining archival
retention and recitation. In such a case, constraints continue to be
enforced, the stated rationale continues to be reproduced, but the
institution has lost the ability to apply the original reasoning to novel
circumstances. This is the characteristic state of an institution undergoing
silent projection collapse: formally intact, locally coherent, and drifting.
\end{observation}

The reason retention principle has several immediate consequences.

It explains why institutional memory is not reducible to document management.
Documents are necessary but not sufficient. What must also be preserved is
the interpretive community capable of engaging those documents --- the people,
practices, and relationships through which the reasoning embedded in constraints
is transmitted across time and membership change.

It explains why apprenticeship, mentorship, and extended professional formation
persist in domains where orientation matters. The explicit content of legal
training, medical education, or scientific apprenticeship can largely be
transmitted through documents. What cannot be transmitted through documents
alone is the capacity to recognize which constraints apply in which situations,
how to weigh competing considerations, and when the formal rule fails to
capture the underlying reason. That capacity is transmitted through extended
engagement with practitioners who already possess it.

It explains why institutional turnover is more dangerous than it appears.
When experienced members leave and are replaced by newcomers whose training
was primarily formal, the institution may retain all its rules while losing
the interpretive capacity required to apply them correctly. The loss is
invisible in the immediate term because the rules continue to be followed.
It becomes visible only when novel situations arise that the rules do not
clearly cover, at which point the institution must either defer to the rule's
surface form or reconstruct a rationale it no longer carries.

It explains why reforms that replace existing constraints with formally
superior alternatives sometimes destroy more than they create. If the
existing constraint embodied an accumulated understanding that was not
fully articulated in its stated rationale, replacing the constraint
eliminates that understanding along with the rule. The new constraint
may be better by every measurable criterion and still leave the institution
less oriented than before, because it has not yet accumulated the interpretive
history that would allow it to be applied intelligently.

\subsection*{Institutional Forgetting as an Economic Pressure}

Reason decay does not occur by accident. It occurs because the activities
required to prevent it are expensive and their returns are invisible until
the capacity is lost.

Apprenticeship is expensive. Extended mentorship is expensive. Deliberation
is expensive. The maintenance of communities of practice, the writing of
case reasoning rather than rule summaries, the study of institutional history
including past failures --- all of these consume time and attention that could
be directed toward measurable outputs. An institution under pressure to
increase throughput discovers that the easiest reductions are precisely the
orientation-preserving activities that produce no immediate measurable return.

The economic structure of this trade-off is asymmetric. The cost of reducing
orientation infrastructure is immediate and measurable: fewer hours in
apprenticeship, faster onboarding, more efficient credentialing. The cost
of that reduction is deferred and difficult to attribute: the loss of
reconstructive capacity that accumulates quietly across years of personnel
change. By the time the cost appears, in the form of poor judgment on novel
situations or inability to revise constraints intelligently, the institution
has long since forgotten the investment decisions that caused it.

The tragedy of institutional forgetting is therefore not that institutions
choose to drift. It is that they repeatedly choose local efficiency over
orientation capacity without recognizing that those choices are in tension.
The institution appears more efficient precisely while becoming less capable
of preserving its purpose.

% -----------------------------------------------------------
\section{Ritualization}

Between full reconstruction and complete abandonment, a constraint can enter
a third state that is probably more common than either: it can become ritualized.

A ritualized constraint is followed but no longer understood. The rule
persists. The behavior persists. What has been lost is the capacity to apply
the underlying reasoning to novel situations, to recognize when the rule
should be revised, and to distinguish cases where the constraint is genuinely
relevant from cases where it is being applied mechanically. The institution
continues to enforce the constraint and to produce the stated rationale on
demand. But the rationale is recited rather than reconstructed, and the
application is formal rather than intelligent.

\begin{observation}
A constraint exists in one of three states: reasoned, when the institution
can reconstruct the argument that justifies it; ritualized, when the
institution follows the constraint and can recite a rationale but cannot
apply the underlying reasoning to novel situations; or abandoned, when the
constraint is neither followed nor remembered. Most institutional theory
attends to the transition between reasoned and abandoned while neglecting
the intermediate state. Ritualization is likely the most common condition
of inherited institutional constraints.
\end{observation}

Ritualization is not purely negative. A ritualized constraint preserves
behavior even after the understanding behind it has been lost. In stable
environments, this may be sufficient. The institution continues to function
correctly for the class of situations the original designers anticipated,
because the behavior the rule prescribes remains appropriate even when the
reason for it is no longer accessible.

The cost of ritualization appears when circumstances change. A novel situation
arises that the original rule did not clearly address. An institution with
reconstructive capacity can apply the underlying reasoning to the new case,
adapt the rule intelligently, or recognize when the constraint has been
made obsolete by changed conditions. An institution that can only recite
a rationale cannot do any of these things. It can apply the rule's surface
form to the new situation, which may produce absurd results. It can make an
exception, which it cannot justify except by appeal to intuition. Or it can
refuse to act because the rule does not clearly apply, which may produce
paralysis.

Ritualization therefore acts as a slow accumulation of institutional
fragility. The institution appears functional in normal conditions and becomes
brittle precisely in the situations that most require intelligent judgment.
The brittleness is not visible until the novel situation occurs. This is
why institutional failure so often appears to originate in exceptional
circumstances: the exceptional circumstance did not cause the failure but
rather revealed the fragility that ritualization had been accumulating.

The transition from reasoned to ritualized is the most common form of reason
decay and the hardest to detect, because it produces no immediate change in
institutional behavior. The rule is still followed. The rationale is still
stated. Only the capacity for reconstruction has disappeared, and that
capacity leaves no visible trace when it goes.

% -----------------------------------------------------------
\section{Institutional Projection Collapse}

The machinery of projection collapse from the previous essay carries over
to the institutional setting with one critical modification.

In individual inquiry, projection collapse occurs when the agent's operative
goal at each step is a projection of the original goal onto the locally
visible context. The projection discards aspects of the original goal that
are not visible in the current context. When the context is impoverished,
optimization against the projection diverges from the original aim. But in
individual inquiry, the original goal remains at least in principle accessible
to the agent. The agent can, through deliberate reflection, compare current
actions against the original objective and detect divergence.

In institutional projection collapse, this comparison operation itself becomes
unreliable. The institution's representation of its original objective is not
stored in any single mind. It is distributed across the membership, embedded
in practices, encoded imperfectly in documents, and partially carried in the
interpretive community whose capacity for reconstruction has been established
above. As that representation degrades --- through turnover, through the
accumulation of local adaptations that gradually displace original meanings,
through the loss of practitioners who carried interpretive capacity informally
--- the institution's ability to compare its current trajectory against its
founding objective diminishes.

\begin{defn}[Institutional Projection Collapse]
An institution undergoes projection collapse when its operative objectives
at each decision cycle are projections of the founding objective onto the
currently visible institutional context, and when the cumulative displacement
between current operative objectives and the founding objective grows across
decision cycles even as each individual transition appears locally reasonable.
Institutional projection collapse is advanced when the institution has lost
sufficient reconstructive capacity that it cannot detect the displacement
by internal comparison.
\end{defn}

The institutional version of projection collapse has a feature that the
individual version does not: it can become self-sealing.

In individual inquiry, even an agent who has drifted substantially from their
original objective retains the cognitive capacity to notice the drift if
prompted to compare current activity against the original aim. The original
aim was present in the same mind that is now engaged in drifted activity.
The comparison is difficult but possible.

In an institution that has lost reconstructive capacity, the comparison is
not merely difficult. The institution lacks the representation of the original
objective against which comparison would be made. Every internal metric
is consistent with the current operative objectives. Every procedure is
internally coherent. Every incentive is aligned with the current projection.
The institution measures itself against its present configuration and finds
that configuration satisfactory. Drift is not detectable from within because
the mechanisms required to detect it --- access to the original objective,
the capacity to reconstruct it, the interpretive community that could
recognize divergence --- have themselves degraded.

\begin{principle}[Invisibility of Advanced Drift]
Projection collapse is structurally invisible to a sufficiently drifted
institution. An institution cannot directly observe the disappearance of an
objective it no longer represents.
\end{principle}

This proposition explains why institutional failures so often appear sudden
to external observers. The failure did not occur suddenly. It accumulated
across many locally reasonable decisions over extended time. But internal
metrics showed continuous improvement, or at least acceptable performance,
throughout the accumulation. The moment of visible failure is not the moment
the institution began to drift. It is the moment external circumstances made
the drift impossible to ignore --- when the gap between current trajectory
and original objective became large enough to produce unmistakable results,
or when an external frame of reference made the drift legible for the first
time.

The self-sealing character of advanced institutional drift also explains
why reform is so difficult. A reform that originates entirely inside a
drifted institution is generated from within the drifted context. It will
tend to optimize against the current operative objectives rather than against
the founding objective. A reform that originates from outside may correctly
identify the drift but lacks the interpretive capacity to reconstruct what
was lost, and so may replace one projection with another rather than
restoring the original orientation.

% -----------------------------------------------------------
\section{The Hierarchy of Institutional Orientation}

Just as orientation in individual inquiry comprises several distinct
capacities that scale differently, institutional orientation comprises
several distinct functions that are differently susceptible to degradation.

The first is \emph{evaluation}: the capacity to measure performance against
explicit criteria. Evaluation is the most tractable level of institutional
orientation because it can often be formalized, automated, and distributed.
Audits, inspections, rankings, and metrics are all instruments of institutional
evaluation. This level of orientation is the most resilient to turnover and
the most amenable to scaling. It is also the most susceptible to substituting
measurable proxies for genuine objectives --- the metric replaces the goal
it was designed to approximate.

The second is \emph{goal maintenance}: the active preservation of the
institution's operative objectives across time and membership change. Goal
maintenance requires more than formal statement of mission. It requires
ongoing engagement with the question of whether current activities serve
the stated mission, and the capacity to resist local adaptations that
gradually displace the mission without formally revising it. Goal maintenance
is more fragile than evaluation because it cannot be fully formalized. It
depends on the interpretive community that carries the reconstructive capacity
described above.

The third is \emph{trajectory assessment}: the capacity to evaluate whether
the institution's current direction, not merely its current state, remains
aligned with its objectives. Trajectory assessment requires comparing not
only where the institution is but where it is going and whether that direction
remains appropriate. This is the level of orientation most commonly missing
in institutional practice. Evaluation asks whether current performance is
acceptable. Trajectory assessment asks whether the current vector will remain
appropriate as circumstances develop. Most institutions conduct evaluation
routinely and trajectory assessment almost never.

The fourth is \emph{goal revision}: the capacity to recognize that the
founding objective was itself mistaken or incomplete, and to revise it without
merely surrendering to drift. Goal revision is the rarest and most difficult
form of institutional orientation. It requires the ability to distinguish
genuine revision --- the principled updating of objectives in response to
new understanding --- from projection collapse masquerading as revision.
Most institutions cannot make this distinction reliably. Constitutional
moments, paradigm shifts, and genuine institutional renewals are the
exceptional cases where goal revision occurs. They are exceptional precisely
because the capacity required for genuine goal revision is so demanding.

\begin{observation}
Most institutions are competent at evaluation, moderate at goal maintenance,
poor at trajectory assessment, and almost incapable of genuine goal revision
except under the pressure of crisis. This is not a contingent feature of
particular institutions. It reflects the relative difficulty of maintaining
each level of orientation under normal institutional pressures.
\end{observation}

The hierarchy also explains why the substitution of lower-level for higher-level
orientation is so common and so destructive. An institution that has lost
trajectory assessment replaces it with more intensive evaluation. More metrics,
more audits, more rankings. These can produce locally satisfying results ---
the metrics improve, the audits pass, the rankings rise --- while the
institution continues to drift, because the intensified evaluation is conducted
against current operative objectives rather than against the founding aim.
The institution becomes more rigorously oriented to its current projection
and less capable of recognizing the gap between that projection and its
original purpose.

% -----------------------------------------------------------
\section{Repair as Reconstruction}

In the previous essay, repair was characterized as the mechanism by which a
living system maintains navigability through time. Repair preserves the
conditions that keep future states reachable. Applied to individual inquiry,
this means preserving structural properties that allow future modification:
keeping interfaces clean, maintaining conceptual coherence, preventing
commitments from accumulating to the point where they foreclose later revision.

In the institutional setting, repair takes a different form. The structural
property that requires preservation is not navigability in a space of actions
but the reconstructive capacity that allows the institution to remain oriented
to its founding objectives. Institutional repair is the work of maintaining
and transmitting the reasons, not merely the rules.

\begin{principle}[Institutional Repair]
Institutional repair is the practice of maintaining reconstructive capacity
across time and membership change. It requires preserving not only the
constraints an institution operates under but the arguments that made those
constraints appropriate, in a form that allows new members to engage them
rather than merely inherit them.
\end{principle}

Institutional repair is not the same as institutional maintenance, which
primarily concerns the preservation of existing structures. It is also not
the same as institutional reform, which primarily concerns the revision of
existing structures. Repair is the work that makes both maintenance and reform
possible without projection collapse: the ongoing transmission of reasons that
allows inherited constraints to be understood, applied intelligently, and
revised when appropriate rather than followed ritualistically or abandoned
unintelligibly.

The activities through which institutional repair is accomplished are not
glamorous and are not easily recognized as productive. They include: extended
apprenticeship and mentorship that transmits interpretive capacity rather than
only formal knowledge; deliberate engagement with institutional history,
including the failures and the reasoning behind rules that exist because
something went wrong; the maintenance of communities of practice in which
the reasons behind constraints are actively discussed rather than silently
assumed; the writing and transmission of case reasoning rather than merely
rule statements; and the preservation of dissent and minority views that
carry alternative framings of the institution's objectives.

These activities are costly relative to their immediate returns. They consume
time and attention that could be directed toward evaluation, production, or
expansion. Under pressure to increase search velocity, they are typically
the first activities to be reduced. An institution under pressure to act
faster will reduce the consultation required before action. It will shorten
apprenticeship, compress mentorship, replace case reasoning with rule lookup,
and substitute formal credentials for demonstrated capacity. Each of these
reductions appears locally rational. The accumulated result is the degradation
of the reconstructive capacity on which long-term institutional orientation
depends.

\begin{observation}
Compression of apprenticeship, replacement of case reasoning by procedure
lookup, and substitution of credentialing for demonstrated interpretive
competence frequently appear as efficiency improvements while functioning
as mechanisms of reason decay. The institution appears more productive
precisely while becoming less capable of preserving its purpose.
\end{observation}

% -----------------------------------------------------------
\section{Governance as Orientation}

Governance is the institutional form of orientation. It is the set of
mechanisms by which a collective preserves coherent objectives across time,
scale, and membership change --- not merely the formal rules by which decisions
are made, but the practices through which the reasons behind those rules
are maintained and transmitted.

This characterization differs from the standard economic account of governance,
which treats it primarily as a coordination mechanism. Coordination is real
and important. But governance as coordination is governance at the evaluation
level: it concerns whether decisions can be made, implemented, and enforced.
Governance as orientation concerns whether the decisions being made remain
coherent with the objectives the governance structure was designed to serve.

\begin{defn}[Governance as Orientation]
Governance, in the sense intended here, is the totality of mechanisms through
which a collective maintains reconstructive capacity across the four levels
of institutional orientation: evaluation, goal maintenance, trajectory
assessment, and goal revision. The adequacy of a governance structure is
not only a question of whether it can produce decisions, but whether those
decisions preserve continuity of aim across time.
\end{defn}

On this characterization, many features of governance that appear irrational
from a pure coordination perspective become intelligible as orientation
mechanisms. Deliberative delay forces consultation of accumulated context
before action. Precedent requires that new decisions be related to prior
reasoning rather than generated independently. Constitutional constraint
requires that major departures from founding objectives be achieved through
processes that make the departure explicit and legible. Separation of powers
creates adversarial processes that partially substitute for trajectory
assessment by generating external comparisons. Long terms and tenure protect
the independence of those who must maintain institutional memory against
short-term pressure.

None of these mechanisms are perfectly efficient. All of them impose costs.
All of them can be captured or corrupted. But their inefficiency, understood
as increased contextual depth, is precisely the feature that makes them
function as orientation mechanisms. They slow the institution down at the
point where slowing down preserves the capacity for reason reconstruction
that would otherwise be lost to search velocity pressure.

The deepest failures of governance are therefore not primarily failures of
coordination, incentive alignment, or even competence. They are failures of
orientation. An institution can maintain perfect formal structure, complete
procedural compliance, and continuous high performance on its current metrics
while undergoing advanced projection collapse. The failure appears when
external circumstances make the gap between current trajectory and founding
objective impossible to ignore. At that point, the institution discovers
that its repair capacity has been eroded precisely during the period when
its coordination capacity appeared most effective.

\subsection*{Adversarial Orientation}

Many of the most durable governance mechanisms preserve orientation not through
consensus but through structured disagreement. Peer review, appellate courts,
opposition parties, independent inspectorates, and red teams all share a common
structure: they create an independent representation of the institution's
objective and compare it against current practice. The disagreement is not
a failure of coordination. It is the orientation mechanism.

Formally, adversarial processes maintain multiple simultaneous projections of
the founding objective $G$. Where a single projection $\pi(G, C_t)$ may drift
without any internal mechanism detecting the divergence, multiple independent
projections $\pi(G, C_t^{(1)}), \pi(G, C_t^{(2)}), \ldots$ generated from
different contextual positions will disagree as drift advances. The disagreement
is itself the signal. An institution whose internal projections remain in
perfect agreement may be either genuinely oriented or so thoroughly drifted
that all its projections have converged on the same displaced objective.
Adversarial processes make that ambiguity legible.

This observation reframes a common institutional tension. The cost of adversarial
processes --- delay, conflict, the need to justify decisions to critical
interlocutors, the possibility of being overturned --- is regularly cited as
evidence that such processes are inefficient. The reframing is that these costs
are orientation infrastructure. They are what makes the independent projection
possible. An institution that eliminates adversarial processes to increase
decision velocity is trading orientation capacity for search velocity, often
without recognizing the exchange.

% -----------------------------------------------------------
\section{Formal Interpretation}

The projection collapse model from the previous essay applies directly to
the institutional setting with one structural extension.

Let $G$ denote the founding objective of an institution. At each decision
cycle $t$, the institution's operative objective is a projection
$\pi(G, C_t)$ of $G$ onto the institutional context $C_t$. As before,
projection is lossy: $\pi(G, C_t)$ preserves aspects of $G$ visible in
$C_t$ and discards the rest.

The institutional context $C_t$ differs from the individual context in a
crucial respect. It is not determined solely by the current state of affairs
but by the composition, memory, and interpretive capacity of the institution's
membership at time $t$. Formally:

\[
C_t \;=\; C_t^{\text{state}} \;\cup\; C_t^{\text{memory}} \;\cup\; C_t^{\text{capacity}}
\]

where $C_t^{\text{state}}$ captures currently visible institutional conditions,
$C_t^{\text{memory}}$ captures accessible recorded history and archived
reasons, and $C_t^{\text{capacity}}$ captures the reconstructive capacity
of the current membership --- the ability to engage, interpret, and apply
the recorded history to novel situations.

\begin{defn}[Reason Decay]
An institution undergoes reason decay at rate $r > 0$ when $C_t^{\text{capacity}}$
decreases over time independently of changes in $C_t^{\text{state}}$ and
$C_t^{\text{memory}}$. Reason decay occurs when the institution loses
practitioners who carry interpretive capacity, when apprenticeship is
compressed, or when formal credentials substitute for demonstrated
reconstructive ability.
\end{defn}

Reason decay is the institutional analog of context impoverishment in
individual projection collapse. As $C_t^{\text{capacity}}$ decreases, the
projection $\pi(G, C_t)$ becomes less faithful even if $C_t^{\text{memory}}$
remains rich. The institution may retain every relevant document and still
be unable to apply the reasoning those documents embody.

\begin{proposition}[Reason Decay Dominance]
For fixed memory retention, sufficient decline in reconstructive capacity
produces projection collapse even when institutional archives remain intact.
Formally: there exist institutions for which $C_t^{\text{memory}}$ remains
constant while $C_t^{\text{capacity}} \to 0$, such that the projection
$\pi(G, C_t)$ diverges from $G$ and the divergence is undetectable by
internal comparison.
\end{proposition}

This proposition formalizes the essay's central empirical claim: archives
alone do not preserve orientation. A complete and well-maintained archive
in the hands of a membership that has lost reconstructive capacity is
orientation infrastructure without orientation. The documents exist. The
capacity to make them do orienting work does not.

The Invisibility Proposition follows formally from this structure. Detection
of projection collapse requires a comparison between $G$ and $\pi(G, C_t)$.
This comparison is possible only if the institution maintains a representation
of $G$ with sufficient fidelity to serve as a reference. A representation
of $G$ is maintained with sufficient fidelity only if $C_t^{\text{capacity}}$
is rich enough to reconstruct the original objective from available memory.
When reason decay has sufficiently impoverished $C_t^{\text{capacity}}$,
the institution cannot maintain a faithful representation of $G$, the
comparison operation fails, and drift becomes internally undetectable.

% -----------------------------------------------------------
\section*{Conclusion}

This essay has argued that institutional drift is structurally inevitable
rather than contingent on bad agents or misaligned incentives, because
constraints persist more readily than the reasons that justified them.
The Reason Retention Principle names the condition for institutional
orientation: an institution remains oriented only to the extent that it
can reconstruct the arguments that originally justified its constraints.
Reconstruction, not mere archival preservation, is the operative standard.

From this principle, several consequences follow. Institutional projection
collapse occurs when accumulated drift displaces operative objectives from
founding objectives across many locally reasonable decisions. Advanced
projection collapse becomes self-sealing when the institution has lost
the reconstructive capacity required to detect the drift. Governance is
not primarily a coordination mechanism but an orientation mechanism: the
totality of practices through which a collective maintains reconstructive
capacity across time and membership change. Institutional repair is not
maintenance or reform but the ongoing transmission of reasons that makes
both possible without collapse.

The four levels of institutional orientation --- evaluation, goal maintenance,
trajectory assessment, and goal revision --- scale differently. Evaluation
is the most tractable and the most susceptible to metric substitution.
Goal revision is the rarest and requires the greatest reconstructive capacity.
The characteristic failure mode is not incompetence at any level but the
substitution of lower-level for higher-level orientation: intensified evaluation
replacing trajectory assessment, producing local metric improvement and
continued directional drift.

The central question for any institution is therefore not whether it is
performing well by its current metrics. It is whether it can still reconstruct
the reasons for which those metrics were chosen. An institution that can
answer that question is oriented. An institution that cannot is drifting,
whether or not any internal measure reflects the fact.

\bigskip

This essay is the second in a pair. The first examined continuity of aim
within an individual mind navigating an expanding possibility space. The
present essay examines continuity of aim within a collective navigating
historical time. In both cases, the central problem emerges when the
production of possibilities grows faster than the capacity to preserve
orientation. Search expands the space of available actions. Orientation
preserves the reason actions are chosen. The scarce resource, in an age
of abundant generation, is not production but continuity. Minds require
it across iterations. Institutions require it across generations. The
medium differs; the structure is the same.

% -----------------------------------------------------------
\bigskip
\begin{center}
  \rule{0.4\textwidth}{0.4pt}
\end{center}

